% Part: incompleteness
% Chapter: theories-computability
% Section: introduction

\documentclass[../../../include/open-logic-section]{subfiles}

\begin{document}

\olfileid{inc}{tcp}{int}

\olsection{مقدمة}

\begin{editorial}
ينبغي إعادة كتابة هذا القسم.
\end{editorial}

لدينا ما يأتي:
\begin{enumerate}
\item تعريف معنى أن تكون دالة قابلة للتمثيل في $\Th{Q}$
  (\olref[req][int]{defn:representable-fn})
\item تعريف معنى أن تكون علاقة قابلة للتمثيل في $\Th{Q}$
  (\olref[req][rel]{defn:representing-relations})
\item مبرهنة تقرر أن الدوال القابلة للتمثيل في $\Th{Q}$ هي بالضبط
  الدوال القابلة للحساب
  (\olref[req][int]{thm:representable-iff-comp})
\item مبرهنة تقرر أن العلاقات القابلة للتمثيل في $\Th{Q}$ هي بالضبط
  العلاقات القابلة للحساب
  (\olref[req][rel]{thm:representing-rels})
\end{enumerate}

{\em النظرية} مجموعة جمل مغلقة استنباطيًا؛ أي لها الخاصية القائلة إنه
متى أثبتت $T$ الصيغة $!A$ انتمت $!A$ إلى $T$. ولعل الأفضل أن نفكر في
النظرية بوصفها مجموعة من الجمل مع كل ما تستلزمه هذه الجمل. ومن الآن
فصاعدًا سنستعمل $\Th{Q}$ للإشارة إلى {\em النظرية} المؤلفة من مجموعة
الجمل القابلة للاشتقاق من البديهيات الثماني في
\olref[req][int]{sec}. وتذكّر أننا نستطيع ترميز صيغ $\Th{Q}$ بأعداد؛
فإذا كانت $!A$ صيغة كهذه، دلّت $\Gn{!A}$ على العدد الذي يرمزها. وبعد
تثبيت هذا الترميز يمكننا أن نسأل ما إذا كانت مجموعات مختلفة من الصيغ
قابلة للحساب أم لا.

\end{document}
